\section{Background Knowledge}
\label{sec:background}

\subsection{Gaussian Processes}

Gaussian Processes (GPs) are a non-parametric Bayesian approach to regression that define a prior over functions rather than a finite set of parameters \parencite{10.7551/mitpress/3206.001.0001}. A GP is fully specified by a mean function $m(x)$ and a covariance function (kernel) $k(x, x')$, such that for any finite set of inputs $\{x_i\}_{i=1}^n$, the corresponding function values follow a multivariate Gaussian distribution as equation \ref{eq:gp_prior}:

\begin{equation}
f(x) \sim \mathcal{GP}(m(x), k(x, x')), \quad
\mathbf{f} = [f(x_1), \dots, f(x_n)]^\top \sim \mathcal{N}(\mathbf{m}, \mathbf{K}).
\label{eq:gp_prior}
\end{equation}

The kernel function encodes assumptions about the structure of the function, such as smoothness and periodicity. The Radial Basis Function (RBF) kernel is defined as equation \ref{eq:rbf_kernel}:
\begin{equation}
k(x, x') = \sigma^2 \exp\left(-\frac{\|x - x'\|^2}{2\ell^2}\right),
\label{eq:rbf_kernel}
\end{equation}
where $\ell$ controls the length-scale and $\sigma^2$ controls the variance. Given observed data $\mathcal{D} = \{X, y\}$, the posterior predictive distribution at a new input $x_*$ is Gaussian, with equation \ref{eq:gp_posterior}:
\begin{equation}
p(f_* \mid x_*, X, y) = \mathcal{N}(\mu_*, \sigma_*^2),
\label{eq:gp_posterior}
\end{equation}
where $\mu_*$ and $\sigma_*^2$ quantify both the prediction and its uncertainty, making GPs naturally suited to probabilistic forecasting.

\subsection{Variational Inference and ELBO}

Exact GP inference scales as $\mathcal{O}(n^3)$ due to matrix inversion, which is prohibitive for approximately 7,500 trading days. Sparse variational GPs address this by introducing $m \ll n$ inducing points that compress the dataset into a lower-dimensional representation \parencite{titsias2009variational}. These inducing points, denoted $\mathbf{Z}$ with corresponding function values $\mathbf{u}$, approximate the true posterior via a variational distribution $q(u)$. Training maximises the Evidence Lower Bound (ELBO) in equation \ref{eq:elbo}:
\begin{equation}
\mathcal{L}_{\mathrm{ELBO}} = \mathbb{E}_{q(f)}[\log p(y \mid f)] - \mathrm{KL}(q(u) \| p(u)),
\label{eq:elbo}
\end{equation}
which balances data fit against a KL regularisation term and is closely related to minimising the negative log-likelihood. This framework, implemented via GPyTorch \parencite{gardner2021gpytorchblackboxmatrixmatrixgaussian}, enables scalable GP inference while retaining principled uncertainty estimation.

\subsection{MC Dropout as Approximate Bayesian Inference}

Monte Carlo (MC) Dropout \parencite{gal2016dropout} interprets dropout at test time as an approximate Bayesian inference technique over neural network weights. By keeping dropout layers active during inference, each forward pass samples from an approximate posterior over the network parameters. Given $N$ stochastic forward passes producing predictions $\{\hat{y}^{(i)}\}_{i=1}^N$, the predictive mean and variance are estimated as in Equation~\ref{eq:mc_dropout_bg}:
\begin{equation}
\mathbb{E}[y] \approx \frac{1}{N} \sum_{i=1}^{N} \hat{y}^{(i)}, \quad
\mathrm{Var}(y) \approx \frac{1}{N} \sum_{i=1}^{N} \left(\hat{y}^{(i)} - \bar{y}\right)^2.
\label{eq:mc_dropout_bg}
\end{equation}
This provides a computationally efficient way to estimate epistemic uncertainty without requiring explicit Bayesian neural networks.
